\begin{document}

\history{Date of publication xxxx 00, 0000, date of current version xxxx 00, 0000.}
\doi{10.1109/ACCESS.2026.DOI}

\title{Congestion-Tolerant and Auditable Trust Management for the
Internet of Digital Twins}

\author{\uppercase{Abdulelah Alwabel}\authorrefmark{1}, \IEEEmembership{Member, IEEE}
}

\address[1]{Department of Computer Sciences,
    College of Computer Engineering and Sciences,
    Prince Sattam bin Abdulaziz University,
    Al-Kharj, Saudi Arabia}

\tfootnote{This work was supported by Prince Sattam bin Abdulaziz
University under Project PSAU/2026/R/1448.}

\markboth
{A. Alwabel: Congestion-Tolerant and Auditable Trust Management for IoDT}
{A. Alwabel: Congestion-Tolerant and Auditable Trust Management for IoDT}

\corresp{Corresponding author: Abdulelah Alwabel (e-mail: a.alwabel@psau.edu.sa)}

\begin{abstract}
The Internet of Digital Twins (IoDT) trust model mirrors every IoT node
with a cloud-resident twin and isolates a node when any of four binary
factors (energy, packet delivery ratio, traceability and twin
notification) fails.  This rule is fragile under congestion, trusts a
single twin, and keeps verdicts in a mutable database.  We propose
three extensions that keep the factor set: a weighted score smoothed by
an exponential moving average with a per-factor persistence check;
quorum voting over signed votes of replicated twins; and a signed
Merkle audit log kept off the decision path.  In simulation, any rule
requiring two failures removes most false isolations at low
congestion, but only the smoothed score keeps them low as congestion
rises: at $4\%$ per-factor congestion it falsely isolated $1.1\%$ of
honest nodes within 16 rounds, against $93\%$ for the AND and
$14$--$26\%$ for two-failure rules, at the cost of one extra round.
On two public packet-delivery traces it still had the lowest false
isolation in every setting, but real losses are bursty and
concentrated on few nodes, raising its rates to $2.0$--$31\%$.  With
identical evidence at all replicas, voting with two of five replicas
corrupted kept detection at the level of an honest twin,
whereas a corrupted single twin hid its node in every run.  The audit
log rejected every edit by an adversary holding fewer than a quorum of
replica keys.  Without congestion, an attacker that knows the rule can
lie in $37.5\%$ of rounds without being flagged.
\end{abstract}

\begin{keywords}
Audit log, digital twin, fog computing, Internet of Things, majority
voting, reputation, trust management.
\end{keywords}

\titlepgskip=-15pt

\maketitle

% ====================================================================
\section{Introduction}
\label{sec:intro}
% ====================================================================
\IEEEPARstart{T}{he} Internet of Things (IoT) is a network of
interconnected devices in which trust is fundamental to the security,
reliability and scalability of large deployments.  Individual devices
are cheap, constrained and routinely exposed to physical and network
compromise; a survey of IoT vulnerabilities, complemented by an
empirical measurement of exploited IoT devices on the Internet,
documents the practical consequences of this exposure~\cite{neshenko2019demystifying}.
Digital twin (DT) technology, reviewed for industry by Tao
\textit{et al.}~\cite{tao2019digital} and surveyed for the IoT by Minerva, Lee and
Crespi~\cite{minerva2020digital}, mirrors the state of each device so
that its activity can be cross-checked.

Soula \textit{et al.}~\cite{soula2024iodt} turned this idea into the
Internet of Digital Twins (IoDT) trust model.  IoDT assigns one
cloud-resident twin to every physical node and folds four binary
factors --- energy, packet delivery ratio (PDR), traceability and twin
notification --- into a single logical-AND verdict.  In its iFogSim
evaluation the model identifies the two compromised nodes in $92\%$
and $95\%$ of cases and isolates them within $250$--$300$~ms.  Three
limitations remain.

\begin{enumerate}
\item \emph{Fragility under transient congestion.}  Under the AND, a
single factor that drops for one round is enough to flag an honest
node.  The original evaluation already reports a false positive
(node N3) and notes that legitimate nodes can be mistaken for
compromised ones during collisions or high-traffic periods; its
authors suggest adjusting the threshold, or adapting the PDR
threshold under congestion, but do not evaluate such a variant~\cite{soula2024iodt}.
Dynamic-trust schemes~\cite{bao2012dynamic,nitti2014trustworthiness}
instead keep continuous scores that age out transient noise.

\item \emph{A single, corruptible twin and an untested adaptive
adversary.}  The original evaluation considers attackers that send
false notifications and a denial-of-service attacker~\cite{soula2024iodt}.
It does not consider on-off attackers that lie only occasionally to
stay below the detection threshold~\cite{sun2006trust},
nor an adversary that corrupts the twin that watches a node, in which
case nothing else in the pipeline can flag the node.  Threats against
the twin itself are part of the DT threat taxonomy
of~\cite{alcaraz2022digital}.

\item \emph{No tamper-evident record.}  Verdicts are stored in a
central database that can be silently rewritten if the cloud is
compromised.  Blockchain-based IoT and fog frameworks provide tamper
evidence~\cite{dorri2017blockchain,novo2018blockchain,kumar2020blockedge,%
dewanta2021bpt,akram2024tmiodt}, but most of them read from or write
to the chain while a decision is being made, and the latency and
computational overhead of blockchain validation is the reason Soula
\textit{et al.}~\cite{soula2024iodt} give for avoiding blockchain in
real-time IoT.
\end{enumerate}

This paper keeps the IoDT factor set and addresses the three
limitations with three extensions.  Its contributions are:
\begin{itemize}
\item a statement of IoDT trust evaluation as a per-round
classification problem, with metrics that separate per-round false
positives, the operationally relevant \emph{false-isolation} rate,
the mean time to the first false isolation over long operation, and
the cost of quarantine with re-admission (Section~\ref{sec:problem});
\item Extension~I, a weighted continuous score with exponential
moving-average (EMA) smoothing and a per-factor persistence check.
A closed form and a counting argument show that the binary AND errs at
rate $\Theta(q)$ in the per-factor congestion probability $q$, any
rule that needs two failures at $\Theta(q^{2})$, and the smoothed
score, at its default parameters and under independent congestion, at
$\Theta(q^{3})$.  The evaluation therefore compares it with
delay-matched baselines (a debounced AND and a 2-of-4 rule) and
quantifies the lie budget of a rule-aware on-off attacker;
\item Extension~II, quorum voting over signed, round-bound vote
vectors of replicated twins, under which, when the replicas share
their evidence and honest votes arrive in time, up to $k-1$ corrupted
replicas change no verdict, and missing or withheld votes never cause
a new isolation;
\item Extension~III, a signed Merkle audit log of verdicts, votes and
trust state that replicas endorse every $K$ rounds off the decision
path and that can be anchored in a permissioned ledger, with a
tamper-evidence theorem that makes explicit when an external
checkpoint is needed; and
\item a reproducible evaluation, with parameters selected on a
separate tuning seed, covering two calibrations to the original
evaluation, long horizons, stealthy, minimal-footprint and rule-aware
attackers, corrupted, withholding and slow replicas, bursty,
common-mode, cross-factor and uneven congestion, a $100$-sensor
topology, and a trace-driven evaluation on two public packet-delivery
datasets (Intel Lab and FlockLab).  The simulator, configurations, raw results and figure
scripts are provided with the submission.
\end{itemize}

The rest of the paper is organised as follows.
Section~\ref{sec:background} reviews related work,
Section~\ref{sec:problem} formalises the problem, the threat model and
the trust assumptions, Section~\ref{sec:proposedModel} presents the
three extensions with their analysis, Section~\ref{sec:resultsAndDiscussion}
reports the experiments, Section~\ref{sec:discussion} discusses them
and compares the design with related frameworks, and
Section~\ref{sec:conclusion} concludes.

% ====================================================================
\section{Related Work}
\label{sec:background}
% ====================================================================
\emph{Digital twins.}  Tao \textit{et al.}~\cite{tao2019digital}
review the state of the art of digital twins in industry, where a twin
integrates the physical and cyber spaces of a product or process.  Minerva \textit{et
al.}~\cite{minerva2020digital} widen the view to the IoT, and Sleuters
\textit{et al.}~\cite{sleuters2019digital} use twins for
the semi-automated behaviour analysis of large distributed IoT
systems.  Gehrmann
and Gunnarsson~\cite{gehrmann2020digital} propose a twin-based security
architecture for industrial control, Alcaraz and
Lopez~\cite{alcaraz2022digital} catalogue DT security threats in a
layered taxonomy that informs our threat model, and Empl \textit{et
al.}~\cite{empl2022soar4iot} prototype SOAR4IoT, a security orchestration,
automation and response application built on twins.

\emph{Digital twins and blockchain.}  Onwubiko \textit{et
al.}~\cite{onwubiko2023enabling} manage digital-twin documents with
Ethereum smart contracts and IPFS storage: stakeholders approve
document versions by a two-thirds threshold on chain and file hashes
are anchored for traceability; the prototype is evaluated by gas cost
and computes no trust score for devices.  Sasikumar \textit{et
al.}~\cite{sasikumar2023blockchain} connect twins of an industrial IoT
network through a proof-of-authority blockchain whose master nodes
validate the network, and report lower energy consumption and better
data security in simulation.  TMIoDT~\cite{akram2024tmiodt} combines a
digital-twin network, a blockchain-based credit score and a
deep-learning intrusion detector for drone services; it is discussed
in Section~\ref{sec:res-comparison}.

\emph{IoT trust management.}  Bao and Chen~\cite{bao2012dynamic} argue
for trust that adapts over time instead of a verdict computed once,
Nitti \textit{et al.}~\cite{nitti2014trustworthiness} carry the idea
to the social IoT, and EigenTrust~\cite{kamvar2003eigentrust} gives the
canonical recipe for continuous reputation by power iteration.  The
beta reputation system of J{\o}sang and Ismail~\cite{josang2002beta}
updates trust from counts of positive and negative observations with
a forgetting factor.  Sun \textit{et
al.}~\cite{sun2006trust} analyse on-off and bad-mouthing attacks on
trust evaluation, and Chen, Bao and Guo~\cite{chen2016trust} evaluate
trust management for the social IoT under self-promoting,
bad-mouthing, ballot-stuffing and discriminatory attacks.  The
surveys~\cite{yan2014survey,guo2017survey,konsta2023survey} review the
field.  Hussain \textit{et al.}~\cite{hussain2020context} propose a
context-aware trust and reputation model for fog-based IoT in which a
device's context is its role or purpose in the network and the trust
level in $[0,1]$ is updated recursively from previous values, with a
weight of $0.7$ on the current trust level, and multi-source feedback.

\emph{Fog and edge computing.}  Bonomi \textit{et
al.}~\cite{bonomi2012fog} introduce fog computing, Chiang and
Zhang~\cite{chiang2016fog} survey its research opportunities for the
IoT, and Shi \textit{et
al.}~\cite{shi2016edge} push computation further to the edge.

\emph{Blockchain for IoT and fog.}  Dorri \textit{et
al.}~\cite{dorri2017blockchain} build a lightweight chain for the smart
home, and Novo~\cite{novo2018blockchain} uses one for scalable access
management.  Debe \textit{et
al.}~\cite{debe2019iot} add an Ethereum-based reputation system for
public fog nodes.  BlockEdge~\cite{kumar2020blockedge} integrates
blockchain pre-processing into an industrial edge pipeline, and
BPT~\cite{dewanta2021bpt} establishes trust between vehicles before a
fog-computing service through a bid and rating exchange recorded on a
chain.

\emph{Tamper-evident logs and anchoring.}  Tamper evidence does not
require a blockchain.  Haber and Stornetta~\cite{haber1991time} link
time-stamped documents by hashes, Schneier and
Kelsey~\cite{schneier1999secure} make audit logs on untrusted machines
tamper-evident with hash chains, and Crosby and
Wallach~\cite{crosby2009efficient} use history trees to make such logs
efficient to audit.  Certificate
Transparency~\cite{laurie2013rfc6962} publishes signed Merkle tree
heads whose leaf and interior hashes are domain-separated.  The
Keyless Signature Infrastructure of Buldas \textit{et
al.}~\cite{buldas2013ksi} publishes only the root of a hash tree, and
Arbitrum~\cite{kalodner2018arbitrum} executes smart contracts
off-chain while the chain records only hashes of their state.
Extension~III follows this line: it is a signed Merkle log that a
ledger may host but whose guarantee does not depend on one.
Section~\ref{sec:res-comparison} compares the closest trust systems on
the figures they actually report.

\emph{Network-level anomaly detection.}  N-BaIoT~\cite{meidan2018nbaiot}
detects botnet traffic with deep autoencoders.  It observes traffic
rather than behavioural trust factors and is complementary to the
present work.

% ====================================================================
\section{Problem Statement}
\label{sec:problem}
% ====================================================================
For every pair of communicating nodes, the IoDT model
of~\cite{soula2024iodt} evaluates four binary factors.  Energy $E$ is
one when the energy a node consumed lies within the bounds
$[E_{\min},E_{\max}]$; $E_{\max}$ follows from the capacity of the
node's super-capacitor and $E_{\min}$ from the CSMA/CA transmission
and back-off cost, both scaled by the number of packets sent.  The PDR
score is one when the measured PDR reaches a threshold
$\Theta_{\mathrm{PDR}}$ ($90\%$ is the example value
of~\cite{soula2024iodt}).  Traceability $T$ is one when the node's most
recent claimed transaction is present in the trust database, and the
twin-notification score $W$ is one unless the twin of the destination
named in the node's notification denounces the claimed communication
(our reading of the notification check of~\cite{soula2024iodt}).  The
verdict is the logical AND
\begin{equation}
\mathrm{Trust}(i) =
E(i) \wedge \mathrm{PDR}(i) \wedge T(i) \wedge W(i).
\label{eq:binary_and}
\end{equation}

\subsection{System Notation}
\label{subsec:notation}
Let $\mathcal{N}$ be the set of physical nodes and
$\mathcal{N}_c\subset\mathcal{N}$ the unknown set of compromised
nodes.  Time is divided into evaluation rounds $t=1,2,\dots$; in every
round the trust engine computes, for every node $i$, the factor vector
\begin{equation}
\mathbf{x}_t(i) =
\bigl(E_t,P_t,T_t,W_t\bigr)(i)
\in \{0,1\}\times[0,100]\times\{0,1\}^{2}
\label{eq:factor_vector}
\end{equation}
where $P_t$ is the PDR in percent.  A trust-aggregation function $f$
maps the history $\mathbf{x}_{1:t}(i)$ to a verdict
$\hat{y}_t(i)\in\{0,1\}$, with $\hat{y}=0$ meaning ``flagged''.

\subsection{Performance Metrics}
\label{subsec:metrics}
The following metrics are used throughout.
\begin{itemize}
\item The \emph{per-round false-positive rate} over a horizon of $T$
rounds,
\begin{equation}
\mathrm{FPR}=\frac{1}{T\,|\mathcal{N}\setminus\mathcal{N}_c|}
\sum_{i\notin\mathcal{N}_c}\sum_{t=1}^{T}\mathbb{1}[\hat{y}_t(i)=0].
\label{eq:fpr}
\end{equation}
\item The \emph{false-isolation rate}: the share of honest nodes
flagged at least once within the horizon $T$.  IoDT isolates a node as
soon as it is flagged~\cite{soula2024iodt}, so this is the number that
matters for a run.  It grows with $T$; the main tables use $T=16$
rounds, the length of one reference scenario
(Section~\ref{sec:setup}), and Section~\ref{sec:res-long} reports it
for up to $1{,}024$ rounds.
\item The \emph{mean time to first false isolation} (MTTFI) of an
honest node over long operation, estimated as the inverse of the rate
at which flags start on honest nodes; it does not depend on $T$.
Strictly, this is the mean time between flag onsets; when flags are
frequent it slightly overstates the time to the first one (about $8$
against $7$ rounds for the AND at $q=0.04$), which does not change any
comparison below.
\item Under \emph{quarantine with re-admission}, a flagged node is
quarantined and re-admitted after $R$ consecutive rounds in which a
quorum of its twins reports $\tau_t\ge\theta_{r}$ and no persistence
trigger ($R=8$, $\theta_r=0.9$); for the baselines, which have no
$\tau_t$, the condition is $R$ consecutive rounds without a flag.  We report the share of time an honest
node spends in quarantine, the number of false quarantines per node
and hour, and the share of rounds in which an attacker delivers a lie
while it is not quarantined.
\item The \emph{detection rate}: the share of runs in which a
compromised node is flagged at least once within $T$ rounds, together
with the mean \emph{time to detection} (TTD) in rounds, and the
\emph{steady-state per-round detection}, the share of rounds $t\geq3$
in which a compromised node is flagged.
\end{itemize}
\subsection{Threat Model}
\label{subsec:threat}
Following the IoT threat survey of~\cite{neshenko2019demystifying} and
the DT threat taxonomy of~\cite{alcaraz2022digital}, the adversary can
compromise a subset $\mathcal{N}_c$ of physical nodes, corrupt at most
$c$ of the $n$ twin replicas that watch a node, steal replica signing
keys, and rewrite or delete records in the cloud.  Compromised nodes
behave in one of four ways.
\begin{itemize}
\item \emph{Lie (IoDT scenario).}  As in~\cite{soula2024iodt}, the
attacker reports a false destination to its twin, and the twin of the
falsely named destination denounces the claim.  We model a lie as
setting $W=0$ and, because the claimed transaction is not found in the
database, also $T=0$; the second effect is our modelling choice.  The
attacker lies in each round with probability $p\in(0,1]$; $p=1$ is the
original naive attacker and $p<1$ a stealthy one.
\item \emph{Minimal-footprint lie.}  The attacker arranges for the
claimed transaction to be logged, so a lie costs only $W=0$, the
smallest observable deviation.
\item \emph{Rule-aware on-off lie.}  The attacker knows the
aggregation rule and its parameters, observes its own past evidence,
and lies only in rounds in which, absent congestion, no trigger would
fire; once flagged it behaves honestly until it is re-admitted.  This
is the on-off attack of~\cite{sun2006trust} against a known detector.
\item \emph{Corrupted twins.}  The adversary controls $c$ replicas
of the twins that watch the two compromised nodes.  Corrupted replicas
vote \textsc{trust} for the attackers (\emph{shielding}), may also
vote \textsc{flag} for honest nodes (\emph{framing}), or may withhold
their votes.
\end{itemize}
Attackers that never deviate on any factor are indistinguishable from
honest nodes at this layer and are out of scope.

Table~\ref{tab:threats} lists which component the design trusts for
what, and which attacks from the IoT trust literature it covers.

\begin{table}[!t]
\caption{Trust assumptions and threat coverage.}
\label{tab:threats}
\centering
\renewcommand{\arraystretch}{1.1}
\footnotesize
\setlength{\tabcolsep}{2.5pt}
\begin{tabular}{p{0.30\columnwidth} p{0.62\columnwidth}}
\toprule
\textbf{Component / attack} & \textbf{Assumption or coverage} \\
\midrule
Fog relay & Trusted to relay evidence and to enforce isolation; a compromised fog node is out of scope \\
Twin replicas & Up to $k-1$ of $n$ Byzantine: shield, frame or withhold (Sec.~\ref{sec:res-ext2}) \\
Trust engine & Trusted for timeliness; the enforced verdicts are checked against the signed votes at every header commit, so an unsupported verdict is exposed within $K$ rounds plus the append time, but not prevented (Sec.~\ref{sec:ext2}) \\
Log storage, ledger & May delete or rewrite; rewrites are rejected by the verifier (Theorem~\ref{thm:tamper}); deleting the newest windows is detected only against a newer checkpoint \\
Checkpoint channel & Controlled independently of the cloud operator \\
\midrule
Lies, stealthy lies & Covered for frequent lies; liars at $p\le0.25$ are detected slowly (Sec.~\ref{sec:res-ext1}, \ref{sec:res-stealth}) \\
Rule-aware on-off & Partly: a lie budget remains (Sec.~\ref{sec:res-stealth}) \\
Bad-mouthing via $W$ & Not covered: a compromised destination can denounce an honest neighbour; repeated denunciations trigger the persistence check \\
Ballot-stuffing via $W$ & Partly: a colluding destination can hide the $W$ failure; under our lie model the lie still fails $T$, but a liar that also gets its transaction logged is not detected \\
Sybil, whitewashing & Out of scope: identities are provisioned by the operator~\cite{douceur2002sybil,friedman2001social} \\
Denial of service & Detected in~\cite{soula2024iodt} mainly through abnormal energy use (a DoS against a twin); not re-evaluated here \\
\bottomrule
\end{tabular}
\end{table}

